\documentclass[10pt, aspectratio=169]{beamer}
\input{../../_en_preambulo_beamer}
% Comandos y macros estadísticas compartidas para las presentaciones Beamer en inglés (Metropolis)

% Conjuntos numéricos básicos
\newcommand{\R}{\mathbb{R}}
\newcommand{\N}{\mathbb{N}}
\newcommand{\Q}{\mathbb{Q}}
\newcommand{\Z}{\mathbb{Z}}
\newcommand{\set}[1]{\left\{ #1 \right\}}
\newcommand{\abs}[1]{\left|#1\right|}
\newcommand{\norm}[1]{\left\| #1 \right\|}

% Comandos estadísticos y probabilísticos del curso
\newcommand{\Var}{\operatorname{Var}}
\newcommand{\cov}{\operatorname{Cov}}
\newcommand{\corr}{\rho}
\newcommand{\comb}[2]{\begin{pmatrix} #1 \\ #2 \end{pmatrix}}
\newcommand{\s}{\sigma}
\newcommand{\particion}{\mathfrak{P}}
\newcommand{\card}[1]{\#\left( #1 \right)}
\newcommand{\seq}[3]{\set{#1}_{#2}^{#3}}
\newcommand{\E}{\mathbb{E}}
\newcommand{\Prob}{\mathbb{P}}

% Entornos pedagógicos y cajas en inglés académico natural (soportando tanto nombres en inglés como en español)
\newenvironment{discussion}{%
  \begin{alertblock}{Discussion Question}%
}{%
  \end{alertblock}%
}
\newenvironment{discusion}{%
  \begin{alertblock}{Discussion Question}%
}{%
  \end{alertblock}%
}

\newenvironment{reflection}{%
  \begin{alertblock}{Classroom Reflection}%
}{%
  \end{alertblock}%
}
\newenvironment{reflexion}{%
  \begin{alertblock}{Classroom Reflection}%
}{%
  \end{alertblock}%
}

\newenvironment{paradox}{%
  \begin{alertblock}{Exploratory Paradox}%
}{%
  \end{alertblock}%
}
\newenvironment{paradoja}{%
  \begin{alertblock}{Exploratory Paradox}%
}{%
  \end{alertblock}%
}

\newenvironment{motivation}{%
  \begin{exampleblock}{Motivation: Data Science in Practice}%
}{%
  \end{exampleblock}%
}
\newenvironment{motivacion}{%
  \begin{exampleblock}{Motivation: Data Science in Practice}%
}{%
  \end{exampleblock}%
}

\newenvironment{quickchallenge}{%
  \begin{block}{Quick Team Challenge}%
}{%
  \end{block}%
}
\newenvironment{retoclase}{%
  \begin{block}{Quick Team Challenge}%
}{%
  \end{block}%
}


\title[Contingency and Independence]{Section 07.04: Contingency Tables and Independence Tests}
\subtitle{Independence, Homogeneity, and the $Z^2=\chi^2$ Identity}
\author[J. Castillo Colmenares]{Juliho Castillo Colmenares}
\institute{Tecnológico de Monterrey}
\date{\vspace{-1.5cm}}

\begin{document}

% Slide 1: Title
\begin{frame}[plain]
  \titlepage
\end{frame}

% Slide 2: Roadmap
\begin{frame}{Roadmap --- Chapter 07: Hypothesis Testing}
  \footnotesize
  Where are we in the modular development of the chapter? \pause
  \begin{itemize}\itemsep=0.08em
    \item \textbf{07.01} Foundations: $H_0$ vs. $H_1$, Type I and II Errors, and Power \pause
    \item \textbf{07.02} $Z$ and $t$ Tests for One- and Two-Sample Means \pause
    \item \textbf{07.03} Chi-Squared Goodness-of-Fit Tests \pause
    \item \textbf{07.04} Contingency Tables and Independence Tests \emph{(Today --- Chapter Closing)}
  \end{itemize}
  \vspace{-0.05cm}
  \begin{block}{Session Objective}
    Extend the $\chi^2$ statistic from Section 07.03 (one categorical variable) to \textbf{two} simultaneous categorical variables: formally distinguish between independence and homogeneity tests, and demonstrate the exact algebraic connection between the $\chi^2$ test and the two-proportion $Z$-test in $2\times2$ tables.
  \end{block}
\end{frame}

% Slide 3: Motivation
\begin{frame}{From One Variable to Two Categorical Variables}
  \footnotesize
  \begin{motivacion}
    In Section 07.03 we asked whether \emph{one} categorical variable follows a theoretical distribution. Now we ask something different and richer: are two categorical variables \emph{related}? Does gender influence career choice? Does generation influence preferred monetization model?
  \end{motivacion}

  \vspace{0.15cm}
  \pause
  \begin{itemize}\itemsep=0.1cm
    \item \textbf{Same mathematical formula}, $\sum(O-E)^2/E$, in both cases. \pause
    \item \textbf{Crucial difference:} the experimental design --- how the sample was collected --- determines whether it is called ``independence'' or ``homogeneity''.
  \end{itemize}
\end{frame}

% Slide 4: Independence vs Homogeneity
\begin{frame}{Independence vs. Homogeneity: the Conceptual Difference}
  \footnotesize
  \begin{block}{Independence Test}
    \textbf{One single sample} of size $N$ is drawn from a single population, and classified \emph{a posteriori} by two variables $A$ and $B$. No marginal total is fixed in advance (except $N$). $H_0: P(A_i\cap B_j)=P(A_i)P(B_j)$.
  \end{block}
  \pause

  \vspace{0.1cm}
  \begin{alertblock}{Homogeneity Test}
    $c$ \textbf{independent} samples of predetermined sizes ($n_1,\ldots,n_c$, fixed columns by design) are drawn from $c$ distinct subpopulations, measuring a single categorical variable. $H_0: p_{i1}=p_{i2}=\ldots=p_{ic}$ for each category $i$.
  \end{alertblock}
\end{frame}

% Slide 5: Common statistic
\begin{frame}{The Common Statistic for Both Tests}
  \footnotesize
  \begin{block}{Expected Frequencies and $\chi^2$ Statistic}
    $$E_{ij} = \frac{R_i C_j}{N}, \qquad \chi^2 = \sum_{i=1}^r\sum_{j=1}^c \frac{(O_{ij}-E_{ij})^2}{E_{ij}} \ \sim\ \chi^2_{(r-1)(c-1)}$$
  \end{block}
  \pause

  \vspace{0.1cm}
  \begin{alertblock}{Key Point}
    Despite the conceptual difference in sampling design, the calculation formula, the degrees of freedom $(r-1)(c-1)$, and the decision procedure are \textbf{identical} for independence and homogeneity. The distinction matters for \emph{interpretation}, not for computation.
  \end{alertblock}
\end{frame}

% Slide 6: Z^2 = chi^2 identity
\begin{frame}{The Exact $Z^2=\chi^2$ Identity in $2\times2$ Tables}
  \footnotesize
  \begin{block}{Algebraic Derivation}
    For a $2\times2$ table with cells $a,b,c,d$: $\chi^2 = \dfrac{N(ad-bc)^2}{(a+b)(c+d)(a+c)(b+d)}$.
  \end{block}
  \pause

  \vspace{0.1cm}
  \begin{alertblock}{Connection to the Two-Proportion $Z$-Test}
    Since $Z\sim N(0,1) \implies Z^2\sim\chi^2_1$, and the two-proportion $Z$-test (Section 07.02) is structurally the signed square root of the homogeneity statistic in a $2\times2$ table: exactly $Z^2\equiv\chi^2$ holds. Both tests are mathematically identical for the binary case.
  \end{alertblock}
\end{frame}

% Slide 7: Applications
\begin{frame}{Applications in Data Science and Engineering}
  \footnotesize
  \begin{itemize}\itemsep=0.1cm
    \item \textbf{Market segmentation:} does preferred monetization model depend on user generation? \pause
    \item \textbf{Clinical trials:} does infection rate depend on the treatment arm (vaccine vs. placebo)? \pause
    \item \textbf{Algorithmic fairness auditing:} are a model's decisions independent of a protected attribute? \pause
    \item \textbf{Multi-group comparison (Marascuilo):} identify exactly which group pairs differ after a global rejection.
  \end{itemize}
\end{frame}

% Slide 8: Python Lab (Block 1)
\begin{frame}[fragile]{Python Lab: Independence Test}
  \scriptsize
  $4\times3$ table (vehicle type vs. day), verified against \texttt{scipy.stats.chi2\_contingency} (\texttt{07.04\_contingency\_tables.py}):
  {\lstset{basicstyle=\fontsize{3.8pt}{4.5pt}\selectfont\ttfamily,aboveskip=0.05em,belowskip=0.05em}
  \lstinputlisting[firstline=17, lastline=32]{../../code/07_pruebas_hipotesis/07.04_contingency_tables.py}}
\end{frame}

% Slide 9: Python Lab (Block 2)
\begin{frame}[fragile]{Python Lab: Homogeneity Test}
  \scriptsize
  Comparing three independent generational cohorts, same formula different design:
  {\lstset{basicstyle=\fontsize{4pt}{4.8pt}\selectfont\ttfamily,aboveskip=0.1em,belowskip=0.1em}
  \lstinputlisting[firstline=43, lastline=64]{../../code/07_pruebas_hipotesis/07.04_contingency_tables.py}}
\end{frame}

% Slide 10: Python Lab (Block 3)
\begin{frame}[fragile]{Python Lab: $Z^2=\chi^2$ Identity}
  \scriptsize
  Exact numerical verification of the algebraic equivalence in a $2\times2$ table:
  {\lstset{basicstyle=\fontsize{3.8pt}{4.5pt}\selectfont\ttfamily,aboveskip=0.05em,belowskip=0.05em}
  \lstinputlisting[firstline=67, lastline=82]{../../code/07_pruebas_hipotesis/07.04_contingency_tables.py}}
\end{frame}

% Slide 11: Terminal Output
\begin{frame}[fragile]{Python Lab: Terminal Output}
  \scriptsize
  Standard output produced when running the Python lab script:
  \vspace{0.02cm}
  \begin{block}{Standard Console Output}
    \fontsize{4pt}{5pt}\selectfont
    \begin{verbatim}
=== Block 1: Independence Test ===
chi2 = 8.4896, df=6, p-value = 0.2044

=== Block 2: Homogeneity Test (3 Cohorts) ===
chi2 = 56.1485, df=4, p-value = 0.000000

=== Block 3: Z^2 = Chi^2 Identity ===
chi2 = 7.3846, Z^2 = 7.3846, |chi2-Z^2| = 0.0000000000
    \end{verbatim}
  \end{block}
\end{frame}

% Slide 12A: Exercise Level 1 (Statement)
\begin{frame}{In-Class Exercise --- Level 1: Fundamental (1/2)}
  \footnotesize
  \begin{block}{Problem (Statement, Problem 3.9.3)}
    A road-safety audit compares vehicle type vs. day of week in $n=300$ reports ($4\times3$ table). Compute the $E_{ij}$ matrix, the $\chi^2$ statistic, df, and decide at $\alpha=0.05$ against $\chi^2_{0.05,6}=12.59$.
  \end{block}

  \vspace{0.15cm}
  \pause
  \begin{alertblock}{Approach}
    $E_{ij}=R_iC_j/300$; $\nu=(4-1)(3-1)=6$.
  \end{alertblock}
\end{frame}

% Slide 12B: Exercise Level 1 (Resolution)
\begin{frame}{In-Class Exercise --- Level 1: Fundamental (2/2)}
  \scriptsize
  We compute the 12 expected cells and sum the contributions: \pause

  \vspace{0.05cm}
  \begin{block}{Calculation}
    $$\chi^2 = \sum_{i,j}\frac{(O_{ij}-E_{ij})^2}{E_{ij}} \approx 8.49, \qquad \nu=6$$
  \end{block}

  \vspace{0.05cm}
  \pause
  \begin{alertblock}{Interpretation}
    Since $8.49<12.59$, \textbf{we do not reject $H_0$}: there is no statistically significant relationship between vehicle type and the day of the week the accident occurs.
  \end{alertblock}
\end{frame}

% Slide 13A: Exercise Level 2 (Statement)
\begin{frame}{In-Class Exercise --- Level 2: Operational (1/2)}
  \footnotesize
  \begin{block}{Problem (Statement, Problem 3.9.5)}
    HR analyzes job satisfaction vs. department in $n=250$ employees ($4\times3$ table). Compute $\chi^2$ and determine whether significant dependence exists at $\alpha=0.05$ ($\chi^2_{0.05,6}=12.59$).
  \end{block}

  \vspace{0.15cm}
  \pause
  \begin{alertblock}{Strategy}
    \scriptsize
    Identify which cell contributes the most to the global statistic for the executive diagnosis.
  \end{alertblock}
\end{frame}

% Slide 13B: Exercise Level 2 (Resolution)
\begin{frame}{In-Class Exercise --- Level 2: Operational (2/2)}
  \scriptsize
  We sum the 12 cell-by-cell contributions: \pause

  \vspace{0.05cm}
  \begin{block}{Calculation}
    $$\chi^2 \approx 16.13, \qquad \nu=6$$
  \end{block}

  \vspace{0.05cm}
  \pause
  \begin{alertblock}{Interpretation}
    \scriptsize
    Since $16.13>12.59$, \textbf{we reject $H_0$}: significant dependence exists. Production concentrates the greatest risk focus (excess of dissatisfied employees); Administration shows far fewer dissatisfied employees than expected.
  \end{alertblock}
\end{frame}

% Slide 14A: Exercise Level 3 (Statement)
\begin{frame}{In-Class Exercise --- Level 3: Analytical (1/2)}
  \footnotesize
  \begin{block}{Problem --- $Z^2\equiv\chi^2$ Equivalence (Statement, Section 3.9)}
    Clinical trial: Vaccine ($n_1=150$, 12 infected), Placebo ($n_2=150$, 28 infected). Run the $\chi^2$ homogeneity test ($\nu=1$) and the two-proportion $Z$-test, verifying $Z^2\equiv\chi^2$.
  \end{block}

  \vspace{0.1cm}
  \pause
  \begin{alertblock}{Strategy}
    \scriptsize
    Compute the pooled $\bar p$ for the $Z$ standard error; compare the result with $\chi^2$ computed directly on the $2\times2$ table.
  \end{alertblock}
\end{frame}

% Slide 14B: Exercise Level 3 (Resolution)
\begin{frame}{In-Class Exercise --- Level 3: Analytical (2/2)}
  \scriptsize
  We compute both statistics independently: \pause

  \vspace{0.05cm}
  \begin{block}{Calculation}
    \scriptsize
    $$\chi^2 = 7.3846, \qquad \bar p=2/15,\ \text{SE}_p\approx0.03925, \qquad Z=\frac{0.08-0.18667}{0.03925}\approx-2.7175,\ Z^2\approx7.3846$$
  \end{block}

  \vspace{0.05cm}
  \pause
  \begin{alertblock}{Interpretation}
    \scriptsize
    The values match up to rounding: $Z^2\equiv\chi^2$ is an exact algebraic identity, not a numerical coincidence. Both tests are 100\% equivalent for the binary $2\times2$ case.
  \end{alertblock}
\end{frame}

% Slide 15A: Exercise Level 4 (Statement)
\begin{frame}{In-Class Exercise --- Level 4: Challenging (1/2)}
  \footnotesize
  \begin{block}{Problem --- Closed-Form $2\times2$ Formula (Statement, Section 3.9)}
    Algebraically derive, starting from Pearson's general definition, that in any $2\times2$ table $\chi^2 = \dfrac{N(ad-bc)^2}{(a+b)(c+d)(a+c)(b+d)}$, and relate $(ad-bc)$ to the Odds Ratio.
  \end{block}

  \vspace{0.1cm}
  \pause
  \begin{alertblock}{Strategy}
    \scriptsize
    First show that the 4 deviations $O_{ij}-E_{ij}$ share the same magnitude $|ad-bc|/N$, alternating sign in a checkerboard pattern.
  \end{alertblock}
\end{frame}

% Slide 15B: Exercise Level 4 (Resolution)
\begin{frame}{In-Class Exercise --- Level 4: Challenging (2/2)}
  \scriptsize
  We factor out the common determinant from the four cells: \pause

  \vspace{0.05cm}
  \begin{block}{Calculation}
    \scriptsize
    $$O_{11}-E_{11}=\frac{ad-bc}{N}, \qquad \chi^2 = \frac{(ad-bc)^2}{N}\left(\frac{1}{R_1}+\frac{1}{R_2}\right)\left(\frac{1}{C_1}+\frac{1}{C_2}\right) = \frac{N(ad-bc)^2}{(a+b)(c+d)(a+c)(b+d)}$$
  \end{block}

  \vspace{0.05cm}
  \pause
  \begin{alertblock}{Interpretation}
    \scriptsize
    The factor $(ad-bc)=\det(\mathbf M)$ is the determinant of the table; $\det=0 \iff ad=bc \iff \text{OR}=ad/bc=1$ (total independence). As the Odds Ratio moves away from 1, $|ad-bc|$ grows quadratically and pushes $\chi^2$ toward rejection.
  \end{alertblock}
\end{frame}

% Slide 16: Closing
\begin{frame}{Section and Chapter Closing: Contingency and Independence}
  \begin{reflexion}
    With this section we close Chapter 07: we have built the complete logical framework of hypothesis testing --- from formulating $H_0/H_a$ and the $\alpha$-$\beta$ trade-off, through means ($Z$/$t$), to complete categorical models ($\chi^2$ goodness-of-fit and contingency tests). All applied statistical inference rests on these tools.
  \end{reflexion}

  \vspace{0.2cm}
  \pause
  \begin{block}{Key Takeaways}
    \begin{itemize}\itemsep=0.08cm
      \item \textbf{Independence vs. homogeneity:} same formula, different sampling design.
      \item \textbf{Common statistic:} $\chi^2=\sum(O_{ij}-E_{ij})^2/E_{ij}$, $\nu=(r-1)(c-1)$.
      \item \textbf{$Z^2=\chi^2$ identity:} proportion tests and $2\times2$ contingency tests are mathematically identical.
    \end{itemize}
  \end{block}
\end{frame}

% Slide 17: Modular Perspective
\begin{frame}{Modular Perspective --- The Next Challenge}
  \scriptsize
  \begin{retoclase}
    With Chapter 07 complete, we master comparing \textbf{two} groups or variables. Chapter 08 (Design of Experiments) generalizes this logic to simultaneously comparing \textbf{three or more} groups via one-way Analysis of Variance (ANOVA) and the $F$-test, avoiding the Type I error inflation that would arise from multiple pairwise $t$-tests.
  \end{retoclase}

  \vspace{0.08cm}
  \pause
  \begin{alertblock}{Recommended Preparation}
    \begin{itemize}\itemsep=0.06cm
      \item Review the Fisher-Snedecor $F$-test and its interpretation as a ratio of variances (Section 05.05).
      \item Reflect on why running $\binom{k}{2}$ pairwise $t$-tests inflates the family-wise Type I error (FWER).
    \end{itemize}
  \end{alertblock}
\end{frame}

\end{document}
